% =========================================================================
% PHỤ LỤC & MA TRẬN PHÂN PHỐI CHƯƠNG TRÌNH ĐẦU SÁCH
% BỘ ĐỀ KIỂM TRA ĐỊNH KÌ TOÁN 1 THEO TỪNG BÀI HỌC
% =========================================================================

\begin{center}
  {\LARGE\bfseries\color{teal!90!black} PHỤ LỤC \& MA TRẬN ĐỀ KIỂM TRA TOÁN 1}\\[4pt]
  {\small\itshape (Ban hành kèm theo Tuyển tập Đề kiểm tra định kì theo từng bài học -- GDPT 2018)}
\end{center}
\vspace{2mm}

% --- KHUNG THÔNG TIN TỔNG QUAN & MÃ QR TOÀN KHÓA ---
\begin{tcolorbox}[
  enhanced,
  colback=teal!5,
  colframe=teal!80!black,
  arc=3mm,
  boxrule=1pt,
  left=3mm, right=3mm, top=3mm, bottom=3mm
]
  \begin{minipage}[c]{0.74\textwidth}
    {\bfseries\color{teal!90!black}\large \faGraduationCap~CỔNG KHẢO THÍ \& LUYỆN THI TRỰC TUYẾN TOÁN 1}\\[3pt]
    \small
    $\bullet$ \textbf{Hệ thống:} \url{https://loptoanthaylinh.vn/thi-online/toan-1}\\
    $\bullet$ \textbf{Quy mô:} 10 Chủ đề toàn diện -- 41 Bài kiểm tra bám sát SGK.\\
    $\bullet$ \textbf{Cấu trúc đề:} 3 Phần chuẩn (Trắc nghiệm chọn A-B-C-D, Đúng/Sai, Điền đáp số ngắn).\\
    $\bullet$ \textbf{Tương tác:} Tự động chấm điểm, lưu lịch sử học tập, tích hợp video chữa câu khó.
  \end{minipage}
  \hfill
  \begin{minipage}[c]{0.23\textwidth}
    \centering
    \qrcode[height=2.3cm]{https://loptoanthaylinh.vn/thi-online/toan-1}\\[3pt]
    {\scriptsize\bfseries\faQrcode~Quét QR toàn khóa}
  \end{minipage}
\end{tcolorbox}

\vspace{3mm}

% --- BẢNG MA TRẬN 41 BÀI KIỂM TRA ---
\noindent{\bfseries\large\color{teal!90!black} Phụ lục 1. Danh mục 10 Chủ đề \& 41 Đề kiểm tra định kì}
\vspace{1.5mm}

\small
\noindent
\begin{tabularx}{\textwidth}{|c|c|X|c|c|c|}
  \hline
  \rowcolor{teal!20}
  \textbf{STT} & \textbf{Chủ đề} & \textbf{Tên bài học kiểm tra} & \textbf{Số câu} & \textbf{Thời gian} & \textbf{Mã thi Online} \\
  \hline
  1 & \multirow{6}{*}{\shortstack{\textbf{Chủ đề 1}\\Các số từ\\0 đến 10}} & Bài 01: Các số 0, 1, 2, 3, 4, 5 & 18 & 35 phút & \texttt{T1-CD1-B01} \\
  \cline{1-1}\cline{3-6}
  2 & & Bài 02: Các số 6, 7, 8, 9, 10 & 18 & 35 phút & \texttt{T1-CD1-B02} \\
  \cline{1-1}\cline{3-6}
  3 & & Bài 03: Nhiều hơn, ít hơn, bằng nhau & 18 & 35 phút & \texttt{T1-CD1-B03} \\
  \cline{1-1}\cline{3-6}
  4 & & Bài 04: So sánh số & 18 & 35 phút & \texttt{T1-CD1-B04} \\
  \cline{1-1}\cline{3-6}
  5 & & Bài 05: Mấy và mấy & 18 & 35 phút & \texttt{T1-CD1-B05} \\
  \cline{1-1}\cline{3-6}
  6 & & Bài 06: Luyện tập chung & 18 & 35 phút & \texttt{T1-CD1-B06} \\
  \hline
  7 & \multirow{3}{*}{\shortstack{\textbf{Chủ đề 2}\\Hình phẳng}} & Bài 07: Hình vuông, hình tròn, tam giác, chữ nhật & 18 & 35 phút & \texttt{T1-CD2-B07} \\
  \cline{1-1}\cline{3-6}
  8 & & Bài 08: Thực hành lắp ghép, xếp hình & 18 & 35 phút & \texttt{T1-CD2-B08} \\
  \cline{1-1}\cline{3-6}
  9 & & Bài 09: Luyện tập chung & 18 & 35 phút & \texttt{T1-CD2-B09} \\
  \hline
  10 & \multirow{4}{*}{\shortstack{\textbf{Chủ đề 3}\\Cộng, trừ\\trong PV 10}} & Bài 10: Phép cộng trong phạm vi 10 & 18 & 35 phút & \texttt{T1-CD3-B10} \\
  \cline{1-1}\cline{3-6}
  11 & & Bài 11: Phép trừ trong phạm vi 10 & 18 & 35 phút & \texttt{T1-CD3-B11} \\
  \cline{1-1}\cline{3-6}
  12 & & Bài 12: Bảng cộng, bảng trừ trong phạm vi 10 & 18 & 35 phút & \texttt{T1-CD3-B12} \\
  \cline{1-1}\cline{3-6}
  13 & & Bài 13: Luyện tập chung & 18 & 35 phút & \texttt{T1-CD3-B13} \\
  \hline
  14 & \multirow{3}{*}{\shortstack{\textbf{Chủ đề 4}\\Hình khối}} & Bài 14: Khối lập phương, khối hộp chữ nhật & 18 & 35 phút & \texttt{T1-CD4-B14} \\
  \cline{1-1}\cline{3-6}
  15 & & Bài 15: V vị trí, định hướng trong không gian & 18 & 35 phút & \texttt{T1-CD4-B15} \\
  \cline{1-1}\cline{3-6}
  16 & & Bài 16: Luyện tập chung & 18 & 35 phút & \texttt{T1-CD4-B16} \\
  \hline
  17 & \multirow{4}{*}{\shortstack{\textbf{Chủ đề 5}\\Ôn tập HK1}} & Bài 17: Ôn tập các số trong phạm vi 10 & 18 & 35 phút & \texttt{T1-CD5-B17} \\
  \cline{1-1}\cline{3-6}
  18 & & Bài 18: Ôn tập phép cộng, phép trừ phạm vi 10 & 18 & 35 phút & \texttt{T1-CD5-B18} \\
  \cline{1-1}\cline{3-6}
  19 & & Bài 19: Ôn tập hình học & 18 & 35 phút & \texttt{T1-CD5-B19} \\
  \cline{1-1}\cline{3-6}
  20 & & Bài 20: Ôn tập chung học kì 1 & 18 & 35 phút & \texttt{T1-CD5-B20} \\
  \hline
\end{tabularx}

\newpage

\noindent
\begin{tabularx}{\textwidth}{|c|c|X|c|c|c|}
  \hline
  \rowcolor{teal!20}
  \textbf{STT} & \textbf{Chủ đề} & \textbf{Tên bài học kiểm tra} & \textbf{Số câu} & \textbf{Thời gian} & \textbf{Mã thi Online} \\
  \hline
  21 & \multirow{4}{*}{\shortstack{\textbf{Chủ đề 6}\\Các số\\đến 100}} & Bài 21: Số có hai chữ số & 18 & 35 phút & \texttt{T1-CD6-B21} \\
  \cline{1-1}\cline{3-6}
  22 & & Bài 22: So sánh số có hai chữ số & 18 & 35 phút & \texttt{T1-CD6-B22} \\
  \cline{1-1}\cline{3-6}
  23 & & Bài 23: Bảng các số từ 1 đến 100 & 18 & 35 phút & \texttt{T1-CD6-B23} \\
  \cline{1-1}\cline{3-6}
  24 & & Bài 24: Luyện tập chung & 18 & 35 phút & \texttt{T1-CD6-B24} \\
  \hline
  25 & \multirow{4}{*}{\shortstack{\textbf{Chủ đề 7}\\Độ dài \&\\Đo độ dài}} & Bài 25: Dài hơn, ngắn hơn & 18 & 35 phút & \texttt{T1-CD7-B25} \\
  \cline{1-1}\cline{3-6}
  26 & & Bài 26: Đơn vị đo độ dài & 18 & 35 phút & \texttt{T1-CD7-B26} \\
  \cline{1-1}\cline{3-6}
  27 & & Bài 27: Thực hành ước lượng và đo độ dài & 18 & 35 phút & \texttt{T1-CD7-B27} \\
  \cline{1-1}\cline{3-6}
  28 & & Bài 28: Luyện tập chung & 18 & 35 phút & \texttt{T1-CD7-B28} \\
  \hline
  29 & \multirow{5}{*}{\shortstack{\textbf{Chủ đề 8}\\Cộng, trừ\\PV 100}} & Bài 29: Phép cộng số có hai chữ số với số có một chữ số & 18 & 35 phút & \texttt{T1-CD8-B29} \\
  \cline{1-1}\cline{3-6}
  30 & & Bài 30: Phép cộng số có hai chữ số với số có hai chữ số & 18 & 35 phút & \texttt{T1-CD8-B30} \\
  \cline{1-1}\cline{3-6}
  31 & & Bài 31: Phép trừ số có hai chữ số cho số có một chữ số & 18 & 35 phút & \texttt{T1-CD8-B31} \\
  \cline{1-1}\cline{3-6}
  32 & & Bài 32: Phép trừ số có hai chữ số cho số có hai chữ số & 18 & 35 phút & \texttt{T1-CD8-B32} \\
  \cline{1-1}\cline{3-6}
  33 & & Bài 33: Luyện tập chung & 18 & 35 phút & \texttt{T1-CD8-B33} \\
  \hline
  34 & \multirow{4}{*}{\shortstack{\textbf{Chủ đề 9}\\Thời gian\\Giờ \& Lịch}} & Bài 34: Xem giờ đúng trên đồng hồ & 18 & 35 phút & \texttt{T1-CD9-B34} \\
  \cline{1-1}\cline{3-6}
  35 & & Bài 35: Các ngày trong tuần & 18 & 35 phút & \texttt{T1-CD9-B35} \\
  \cline{1-1}\cline{3-6}
  36 & & Bài 36: Thực hành xem lịch và giờ & 18 & 35 phút & \texttt{T1-CD9-B36} \\
  \cline{1-1}\cline{3-6}
  37 & & Bài 37: Luyện tập chung & 18 & 35 phút & \texttt{T1-CD9-B37} \\
  \hline
  38 & \multirow{4}{*}{\shortstack{\textbf{Chủ đề 10}\\Ôn tập\\cuối năm}} & Bài 38: Ôn tập các số và phép tính phạm vi 10 & 18 & 35 phút & \texttt{T1-CD10-B38} \\
  \cline{1-1}\cline{3-6}
  39 & & Bài 39: Ôn tập các số và phép tính phạm vi 100 & 18 & 35 phút & \texttt{T1-CD10-B39} \\
  \cline{1-1}\cline{3-6}
  40 & & Bài 40: Ôn tập hình học và đo lường & 18 & 35 phút & \texttt{T1-CD10-B40} \\
  \cline{1-1}\cline{3-6}
  41 & & Bài 41: Ôn tập chung cuối năm & 18 & 35 phút & \texttt{T1-CD10-B41} \\
  \hline
\end{tabularx}

\vspace{4mm}

% --- PHỤ LỤC 2: HƯỚNG DẪN QUÉT QR VÀ QUY CHẾ ĐÁNH GIÁ ---
\noindent{\bfseries\large\color{teal!90!black} Phụ lục 2. Hướng dẫn quét QR thi online \& Thang điểm đánh giá}
\vspace{1.5mm}

\begin{tcolorbox}[
  enhanced,
  colback=white,
  colframe=teal!70!black,
  arc=2.5mm,
  boxrule=0.8pt,
  left=3mm, right=3mm, top=2.5mm, bottom=2.5mm
]
  \begin{enumerate}[\bfseries 1.]
    \item \textbf{Cách quét mã QR}:
      \begin{itemize}
        \item Mở ứng dụng \textbf{Camera} hoặc \textbf{Zalo} trên điện thoại/máy tính bảng.
        \item Hướng camera về phía mã QR ở góc trên bên phải mỗi đề thi.
        \item Nhấn vào thông báo liên kết hiện lên để truy cập ngay vào bài kiểm tra online tương ứng.
      \end{itemize}
    \item \textbf{Cấu trúc đề kiểm tra (Tổng 10.0 điểm)}:
      \begin{itemize}
        \item \textbf{Phần 1 (6.0 điểm)}: 12 câu trắc nghiệm 4 lựa chọn (0.5 điểm/câu).
        \item \textbf{Phần 2 (2.0 điểm)}: Trắc nghiệm Đúng / Sai (Kiểm tra tư duy phản biện và khả năng nhận biết).
        \item \textbf{Phần 3 (2.0 điểm)}: Câu hỏi trả lời ngắn / điền số (Kiểm tra kỹ năng tính toán thực tế).
      \end{itemize}
    \item \textbf{Xếp loại năng lực}:
      \begin{itemize}
        \item \textbf{Hoàn thành xuất sắc}: $9.0 - 10.0$ điểm (Thưởng $+5$ sao tích lũy).
        \item \textbf{Hoàn thành tốt}: $7.0 - 8.5$ điểm (Thưởng $+3$ sao tích lũy).
        \item \textbf{Cần cố gắng \& xem lại video bài giảng}: Dưới $7.0$ điểm.
      \end{itemize}
  \end{enumerate}
\end{tcolorbox}
